% ============================================================================
% Section 1.2 — Recognizing Proportional Relationships
% CCSS 7.RP.A.2.A
% ============================================================================
\section{Recognizing Proportional Relationships}

\begin{stepGoal}
\begin{itemize}
    \item Decide whether two quantities are proportional using a table
    \item Decide whether a graph shows a proportional relationship
\end{itemize}
\end{stepGoal}

\begin{stepVocabBox}
    \vocabItem{Proportional}{Two quantities whose ratio stays the \textbf{same} for every pair.}
    \vocabItem{Origin}{The point $(0, 0)$ on a coordinate plane.}
\end{stepVocabBox}

\begin{stepsCard}[How to Tell If a Relationship Is Proportional]
    \stepItem{For each row in the table, divide $y \div x$ to find the ratio.}
    \stepItem{Check: are \textbf{all} the ratios the \textbf{same}?}
    \stepItem{If yes $\to$ \textbf{proportional}. If any ratio is different $\to$ \textbf{not proportional}.}
    \stepItem{On a graph: the points must form a \textbf{straight line through $(0, 0)$}.}
\end{stepsCard}

% --- Worked Example 1 ---
\begin{stepExample}{Is this relationship proportional?
\begin{center}
\begin{tabular}{|c|c|}
\hline $x$ & $y$ \\ \hline $2$ & $\$24$ \\ $5$ & $\$60$ \\ $8$ & $\$96$ \\ \hline
\end{tabular}
\end{center}}
    \stepShow{1} Divide $y \div x$ for each row: $24 \div 2 = 12$, \; $60 \div 5 = 12$, \; $96 \div 8 = 12$

    \stepShow{2} All three ratios equal $12$.

    \stepShow{3} Every ratio is the same $\to$ proportional.
    \stepResult{Yes — proportional at $\$12$ per unit.}
\end{stepExample}

% --- Worked Example 2 ---
\begin{stepExample}{Is this relationship proportional?
\begin{center}
\begin{tabular}{|c|c|}
\hline $x$ & $y$ \\ \hline $2$ & $6$ \\ $4$ & $10$ \\ $6$ & $14$ \\ \hline
\end{tabular}
\end{center}}
    \stepShow{1} Divide: $6 \div 2 = 3$, \; $10 \div 4 = 2.5$, \; $14 \div 6 \approx 2.3$

    \stepShow{2} The ratios ($3$, $2.5$, $2.3$) are NOT all equal.

    \stepShow{3} At least one ratio is different $\to$ not proportional.
    \stepResult{No — the ratios are not constant.}
\end{stepExample}

\watchOut{A straight line alone does NOT mean proportional. The line must also pass through $(0, 0)$. A line that crosses the $y$-axis at $5$ is NOT proportional.}

% ============================================================================
% PRACTICE
% ============================================================================
\newpage
\begin{stepPractice}{Recognizing Proportional Relationships Practice}
    \resetProblems

    \stepPracticeHeader[step1Color]{Divide to Find Each Ratio}
    \prob Find $y \div x$ for each row. Are they equal?
    \begin{center}
    \begin{tabular}{|c|c|}
    \hline $x$ & $y$ \\ \hline $3$ & $9$ \\ $5$ & $15$ \\ $7$ & $21$ \\ \hline
    \end{tabular}
    \end{center}
    \answer{Yes — all ratios equal $3$, so it is proportional.}

    \stepPracticeHeader[step2Color]{Proportional or Not?}
    \prob Is this proportional? Explain.
    \begin{center}
    \begin{tabular}{|c|c|}
    \hline $x$ & $y$ \\ \hline $1$ & $4$ \\ $3$ & $10$ \\ $5$ & $16$ \\ \hline
    \end{tabular}
    \end{center}
    \answer{No — $4 \div 1 = 4$ but $10 \div 3 \approx 3.3$. Ratios differ.}

    \stepPracticeHeader[step3Color]{Graph Test}
    \prob True or false: A straight line through $(0, 0)$ always shows a proportional relationship. \answerBlank[2cm]
    \answerTF{True}

    \prob A graph passes through $(0, 0)$ and $(4, 12)$. Is it proportional? What is the ratio? \answerBlank[2.5cm]
    \answer{Yes — $12 \div 4 = 3$; the ratio is $3$.}

    \wordProblem{A phone plan costs $\$0.05$ per text with no monthly fee. A second plan costs $\$5$/month plus $\$0.02$ per text. Which plan is proportional? Explain.}{~}
    \answer{Plan 1 — no base fee, so cost $\div$ texts is always $\$0.05$. Plan 2 has a fixed $\$5$ fee, so the ratio changes.}
\end{stepPractice}
